
Model choice in MBRL is delicate: we desire effective learning in both low-data regimes (at the beginning) and high-data regimes (in the later stages of the learning process). 
For this reason, Bayesian nonparametric models, such as Gaussian processes (GPs), are often the model of choice in MBRL, especially in low-dimensional problems where data efficiency is critical~\citep{Kocijan2004,Ko2007,Nguyen-Tuong2008,Grancharova2008,Deisenroth2014,Kamthe2018}. 
However, such models introduce additional assumptions on the system, such as the smoothness assumption inherent in GPs with squared-exponential kernels~\citep{Rasmussen2003}.
Parametric function approximators have also been used extensively in MBRL~\citep{Hernandaz1990,Miller1990,Lin1992,Draeger1995}, but were largely supplanted by Bayesian models in recent years.
Methods based on local models, such as guided policy search algorithms~\citep{Levine2016,Finn2015,Chebotar2017}, can efficiently train NN policies, but use time-varying linear models, which only locally model the system dynamics.
Recent improvements in parametric function approximators, such as NNs, suggest that such methods are worth revisiting~\citep{Baranes2013,Fu2016,Punjani2015,Lenz2015,Agrawal2016,Gal2016,Depeweg2016,Williams2017,Nagabandi2017}.
Unlike Gaussian processes, NNs have constant-time inference and tractable training in the large data regime, and have the potential to represent more complex functions, including non-smooth dynamics that are often present in robotics~\citep{Fu2016,Mordatch2016,Nagabandi2017}.
However, most works that use NNs focus on deterministic models, consequently suffering from overfitting in the early stages of learning.
For this reason, our approach is able to achieve even higher data-efficiency than prior deterministic MBRL methods such as~\cite{Nagabandi2017}.

Constructing good Bayesian NN models remains an open problem~\citep{MacKay1992,neal1995bayesian,Osband2016a,Guo2017}, although recent promising work exists on incorporating dropout~\citep{gal2017concrete}, ensembles~\citep{Osband2016,Lakshminarayanan2017}, and $\alpha$-divergence~\citep{hernandez2016black}.
Such probabilistic NNs have previously been used for control, including using dropout~\cite{Gal2016,Higuera2018} and $\alpha$-divergence~\cite{Depeweg2016}.
In contrast to these prior methods, our experiments focus on more complex tasks with challenging dynamics -- including contact discontinuities -- and we compare directly to prior model-based and model-free methods on standard benchmark problems, where our method exhibits asymptotic performance that is comparable to model-free approaches.
